\documentclass[11pt, landscape]{article}

% change the size of the page
\usepackage[
	papersize={20in, 11.5in},
	margin=1in
]{geometry}

\usepackage[table, dvipsnames]{xcolor}
\usepackage{tabularx}
\usepackage{fontspec}
\usepackage{etoolbox}
\usepackage{enumitem}
\usepackage{environ}
\usepackage{xltabular}
\usepackage[most]{tcolorbox}

\setmainfont{PT Sans} 
\usepackage[portuguese]{babel}

% --- Color Palette ---
\definecolor{primaryBlue}{HTML}{2C3E50}
\definecolor{normalGreen}{HTML}{27AE60}
\definecolor{altOrange}{HTML}{E67E22}
\definecolor{excRed}{HTML}{C0392B}
\definecolor{lightGray}{HTML}{F8F9FA}
\definecolor{borderGray}{HTML}{D1D1D1}

% --- Table Styling ---
\renewcommand{\arraystretch}{1.6}
\setlength{\arrayrulewidth}{0.5pt}
\arrayrulecolor{borderGray}

% --- Counters ---
\newcounter{stepcount}
\newcounter{substepcount}

% --- Commands ---

\newcommand{\ucHeader}[5]{
	\noindent {\color{primaryBlue}\Huge \textbf{Caso de Uso: #1}}
	
	\vspace{1.5em}
	\noindent
	\begin{tabularx}{\textwidth}{|l|X|}
		\hline
		\rowcolor{lightGray} \textbf{Descrição} & #2 \\ \hline
		\textbf{Cenários} & #3 \\ \hline
		\rowcolor{lightGray} \textbf{Pré-condições} & #4 \\ \hline
		\textbf{Pós-condições} & #5 \\ \hline
	\end{tabularx}
}

\NewEnviron{flowtable}{
	\noindent
	\begin{xltabular}{\textwidth}{|>{\hsize=1.2\hsize}X|>{\hsize=0.8\hsize}X|>{\hsize=1.4\hsize}X|>{\hsize=0.6\hsize}X|}
		\hline
		\rowcolor{lightGray} 
		\textbf{Passos} & \textbf{Lógica de Negócio} & \textbf{Definição da API} & \textbf{Subsistema} \\ \hline
		\endfirsthead
		
		\hline
		\rowcolor{lightGray} 
		\textbf{Passos} & \textbf{Lógica de Negócio} & \textbf{Definição da API} & \textbf{Subsistema} \\ \hline
		\endhead
		
		\BODY
	\end{xltabular}
}

\newcommand{\step}[4]{
	\stepcounter{stepcount}
	\thestepcount. #1 & #2 & \texttt{#3} & #4 \\ \hline
}

\newcommand{\substep}[4]{
	\stepcounter{substepcount}
	\theparentstep.\thesubstepcount\ #1 & #2 & \texttt{#3} & #4 \\ \hline
}

% lazy step: use it when the operation will not be reflected in the system
\newcommand{\lstep}[1]{
	\step{#1}{-}{-}{-}
}

% lazy sub step: use it when the operation will not be reflected in the system
\newcommand{\lsubstep}[1]{
	\substep{#1}{-}{-}{-}
}

\newcommand{\normalFlowHeader}{
	\setcounter{stepcount}{0}
	
	\vspace{1em}
	
	\noindent
	\begin{tcolorbox}[
		enhanced,
		sharp corners,
		boxrule=0pt,
		frame hidden,
		colback=normalGreen!8!white,
		left=5pt, right=5pt,
		top=8pt, bottom=8pt,
		after skip=0pt,
		]
		{\color{normalGreen}\Large\bfseries Fluxo Normal}
	\end{tcolorbox}
	\nointerlineskip
}

\newcommand{\subFlowHeader}[4]{
	\setcounter{substepcount}{0}
	\def\theparentstep{#3}
	
	\noindent
	\begin{tcolorbox}[
		enhanced,
		sharp corners,
		boxrule=0pt,
		frame hidden,
		colback=#1!8!white,
		left=5pt, right=5pt,
		top=5pt, bottom=5pt,
		after skip=0pt,
		]
		{\color{#1}\Large\bfseries #2} \vspace{0.3em} \\
		\noindent\hspace{2pt}\vrule width 2pt depth 0.8ex height 2.2ex \hspace{0.5em}
		\begin{minipage}{\dimexpr\textwidth-25pt}
			\small \textbf{PASSO #3} \quad \textbullet \quad \textbf{CONDIÇÃO:} \textit{#4}
		\end{minipage}
	\end{tcolorbox}
	\nointerlineskip
}

\begin{document}
	
	\ucHeader{Instalar Painéis}
	{O Actor dirige-se a local de instalação e efectua-a}
	{}
	{Existem instalações a efectuar}
	{Instalação efectuada com sucesso}
	
	\normalFlowHeader
	\begin{flowtable}
		\step{Sistema determina próximas instalações e comunica-as ao Actor}{Obter próximos pedidos de instalação}{getProximosPedidos() : List<Pedido>}{Comercial}
		\lstep{Actor seleciona instalação a realizar}
		\step{Sistema apresenta informação do pedido de instalação}{Consultar informação do pedido de instalação}{getPedido(pedido : String) : Pedido}{Comercial}
		\lstep{Actor indica número do contrato de energia}
		\lstep{Actor indica orientação e potência dos painéis a instalar e consumo típico da casa}
		\step{Sistema valida orientação e informque Actor de que esta é adequada}{Validar orientação e potência dos painéis}{validarPainel(ori : Orientacao, potencia : float, consumo : float) : boolean}{Tecnico}
		\lstep{Actor indica IPs do \textit{meter} e da \textit{box}}
		\step{Sistema regista as informações fornecidas no contrato}{Registar instalação de painéis}{registarInstalação(contrato : String, meterIP : IP, boxIP : IP) : void}{Cliente}
		\step{Actor confirma o bom funcionamento da instalação}{Registar funcionamento da instalação}{registarFuncionamento(contrato : String, estado : boolean) : void}{Cliente}
		\lstep{Sistema encerra o processo}
	\end{flowtable}
	
	\subFlowHeader{altOrange}{Fluxo Alternativo}{4}{Contrato de energia inexistente}
	\begin{flowtable}
		\lsubstep{Actor indica inexistência de contrato de energia}
		\substep{Sistema regista novo contrato de energia}{Registar novo contrato de energia}{criarContrato(pedido : String) : void}{Cliente}
		\lsubstep{Regressa a 5}
	\end{flowtable}

	\newpage

	\subFlowHeader{altOrange}{Fluxo Alternativo}{6}{Produção será excessiva no pico de produção}
	\begin{flowtable}
		\lsubstep{Sistema ao validar orientação conclui que ela será excessiva no pico de produção}
		\substep{Sistema calcula capacidade ideal da bateria a instalar }{Calcular capacidade ideal da bateria}{calcularCapacidade(potencia : float, consumo : float) : int}{Tecnico}
		\substep{Sistema regista possibilidade de instalação de bateria (para contacto posterior) e informa Actor}{Registar contacto para serviço de bateria}{registarContacto(contacto : int, consumo : float) : void}{Comercial}
		\lsubstep{Regressa a 7}
	\end{flowtable}
	
	\subFlowHeader{excRed}{Fluxo de Exceção}{4}{Impossível aceder ao local}
	\begin{flowtable}
		\lsubstep{Actor indica impossibilidade de acesso}
		\substep{Sistema regista impossibilidade de acesso e adia instalação}{Registar impossibilidade de instalação}{registarImpossibilidade(pedido : String, motivo : String) : void}{Comercial}
	\end{flowtable}

	\subFlowHeader{excRed}{Fluxo de Exceção}{9}{Actor indica mau funcionamento da instalação}
	\begin{flowtable}
		\substep{Sistema regista mau funcionamento da instalação}{Registar funcionamento da instalação}{registarFuncionamento(contrato : String, estado : boolean) : void}{Cliente}
	\end{flowtable}
	
\end{document}